\section*{第 \ref{ch:g12:photosynthesis} 章——光合作用}

\begin{solution}{exo:g12:photosynthesis:1}
一种具有双层膜的\omterm{def:g10:cells-common-unit:organelle}{细胞器}，其内充满液态的基质，基质中含有
堆叠排列的\omterm{def:g12:photosynthesis:chloroplast}{类囊体}膜。色素和光能转换蛋白
位于\omterm{def:g12:photosynthesis:chloroplast}{类囊体}膜上；合成糖的酶
位于基质中。
\end{solution}

\begin{solution}{exo:g12:photosynthesis:2}
叶绿素吸收蓝光和红光而反射绿光，这就是我们
所看到的颜色。蓝光和红光最能驱动\omterm{def:g10:cell-metabolism:autotrophy}{光合作用}。
\end{solution}

\begin{solution}{exo:g12:photosynthesis:3}
ATP、还原型 NADP 和氧气；其中氧气是副产物，被释放出来。
\end{solution}

\begin{solution}{exo:g12:photosynthesis:4}
氧气来自被分解的水，而非二氧化碳：给藻类
提供标记重氧的水时，它们释放出含重氧的 $\mathrm{O_2}$；
提供标记二氧化碳时则没有。
\end{solution}

\begin{solution}{exo:g12:photosynthesis:5}
六个 $\mathrm{CO_2}$、18 个 ATP 和 12 个还原型 NADP。
\end{solution}

\begin{solution}{exo:g12:photosynthesis:6}
在 \qty{550}{nm} 处：吸收率约为 8\%，速率约为 25\%。在
\qty{665}{nm} 处：吸收率约为 90\%，速率约为 95\%。绿色灯
提供的是叶片利用最少的光；若温室采用绿光照明，
植物几乎无法生长，所谓节能也就毫无意义。
\end{solution}

\begin{solution}{exo:g12:photosynthesis:7}
细菌向氧气游动，它们聚集在释放氧气的地方，
即藻类进行\omterm{def:g10:cell-metabolism:autotrophy}{光合作用}的地方，也就是光谱中
蓝光和红光的区段。以氧气释放量衡量的\omterm{def:g10:cell-metabolism:autotrophy}{光合作用}，
取决于光的颜色，并遵循色素的
吸收规律。
\end{solution}

\begin{solution}{exo:g12:photosynthesis:8}
氧气是在没有任何碳被固定的情况下释放的：光反应
分解水、传递电子，与碳反应彼此独立，
不需要 $\mathrm{CO_2}$，只需要光、水和
能接收电子的物质。
\end{solution}

\begin{solution}{exo:g12:photosynthesis:9}
没有光，就没有 ATP 和还原型 NADP，因此三碳酸
无法被还原而积累；受体仍能在一段时间内捕获
$\mathrm{CO_2}$，但不能再生，
最终被耗尽。
\end{solution}

\begin{solution}{exo:g12:photosynthesis:10}
光强随距离减弱（大致与距离的平方成反比：距离加倍时，
光强减为四分之一），光反应的速度随之变慢。
在黑暗中不释放氧气——而且植物还会通过
呼吸作用消耗一部分氧气。
\end{solution}

\begin{solution}{exo:g12:photosynthesis:11}
$8/44 \approx \qty{0.18}{mol}$ 的 $\mathrm{CO_2}$：$0.18/6 \approx
\qty{0.030}{mol}$ 葡萄糖，约 \qty{5.5}{g}；\qty{0.18}{mol} 的
$\mathrm{O_2}$，约 \qty{5.8}{g}。
\end{solution}

\begin{solution}{exo:g12:photosynthesis:12}
它本身并不直接利用光，只利用光反应的产物。但
在黑暗中，一旦 ATP 和还原型 NADP 耗尽，
它便在一分钟内停止：它在光照下运行，
由光反应提供能量；“暗”指的是它所需的条件，而非它发生的时间。
\end{solution}

\begin{solution}{exo:g12:photosynthesis:13}
$^{18}$O：一分钟内即出现在释放的 $\mathrm{O_2}$ 中，
此后一小时和一天后仍在其中（部分也出现在呼吸作用生成的水中）。
$^{14}$C：一分钟后出现在循环中的三碳酸和\omterm{def:g10:chemistry-of-life:families}{糖类}中；
一小时后出现在蔗糖和淀粉中；一天后出现在纤维素、
\omterm{def:g10:chemistry-of-life:families}{脂质}、\omterm{def:g10:chemistry-of-life:families}{氨基酸}中，以及经呼吸释放的 $\mathrm{CO_2}$ 中。
\end{solution}

\begin{solution}{exo:g12:photosynthesis:14}
$20/1000 = 2\%$。其余部分被反射（即绿光），
被不适用的波长吸收并转化为热量，用于
使水分蒸发的蒸腾作用，或以热量形式
在光反应本身中散失。
\end{solution}

\begin{solution}{exo:g12:photosynthesis:15}
呼吸作用、燃烧以及岩石和铁的氧化会持续消耗
氧气；若无\omterm{def:g10:cell-metabolism:autotrophy}{光合作用}加以补充，大气中的氧气储量
将在几千年内耗尽。早期大气中没有游离氧，
它只是在光合生物长期产生氧气之后才逐渐积累——
最初的二十亿年
里氧气含量很低。
\end{solution}

\begin{solution}{pb:g12:photosynthesis:1}
\textbf{1.} $8/44 \approx \qty{0.18}{mol}$。

\textbf{2.} $0.18/6 = \qty{0.030}{mol}$，约 \qty{5.5}{g}。

\textbf{3.} \qty{0.18}{mol} 的 $\mathrm{O_2}$，约 \qty{4.4}{L}。

\textbf{4.} 每个 $\mathrm{O_2}$ 需两个水分子：\qty{0.36}{mol}，
约 \qty{6.5}{g}——相当于蒸腾散失的 \qty{600}{mL} 水的百分之一。

\textbf{5.} 氧气来自\omterm{def:g12:photosynthesis:chloroplast}{类囊体}中被分解的水分子：由重氧标记可见，
重氧标记仅在 $\mathrm{O_2}$ 中出现，
且只有水被标记时才出现。

\textbf{6.} 每固定一个 $\mathrm{CO_2}$ 需一轮\omterm{prop:g12:photosynthesis:calvin}{卡尔文循环}：共 \qty{0.18}{mol} 轮。

\textbf{7.} 每轮循环需三个 ATP 和两个还原型 NADP：
共 \qty{0.55}{mol} 的 ATP 与 \qty{0.36}{mol} 的还原型 NADP。

\textbf{8.} 每分解一个水分子产生两个电子，可形成一个还原型 NADP：
\qty{0.36}{mol} 的水，与第 4 题的结果相同，
这是必然的。

\textbf{9.} 氧气的产生会立即停止；还原型 NADP 与 ATP 将不再生成，
并在数秒内耗尽；
糖的合成也将在携带者耗尽后停止。

\textbf{10.} \omterm{prop:g12:photosynthesis:calvin}{卡尔文循环}不再消耗携带者；
无物接收其电子与能量，光反应阶段将在数秒内积压减速——
两个阶段
通过携带者相互耦合。

\textbf{11.} $0.030 \times 2870 \approx \qty{86}{kJ}$ 每小时，即
$86000/3600 \approx \qty{24}{W}$。

\textbf{12.} $24/1000 = 2.4\%$ 占全部日光的比例；$24/450 \approx
5\%$ 占被吸收波长的比例。

\textbf{13.} $8 \times 176 \approx \qty{1400}{kJ}$ 为每摩尔
$\mathrm{CO_2}$ 所对应的光子能量；储存量：$2870/6 \approx \qty{480}{kJ}$。

\textbf{14.} 约三分之一；其余部分以热量形式在
两个阶段连续的能量传递中散失。

\textbf{15.} 葡萄糖净增益 $5.5 - 0.5 = \qty{5}{g}$，每小时。
以光照下的净 $\mathrm{CO_2}$ 吸收量测得；
单独的呼吸作用在黑暗中测量，再加回去，
即得总\omterm{def:g10:cell-metabolism:autotrophy}{光合作用}。

\textbf{16.} $\num{8.5e11}/\num{1e11} \approx 8.5$ 年——
若再计入海洋藻类固定的同等量，约 4 年。

\textbf{17.} 呼吸作用（植物、动物、真菌和细菌的），
连同燃烧，以几乎相同的速率释放二氧化碳；固定与释放
基本平衡，因此在化石燃料燃烧
增加新来源之前，碳储量一直稳定。

\textbf{18.} \qty{1e11}{t} 碳相当于 $\num{8.3e12}\,\unit{mol}$；
同样摩尔数的 $\mathrm{O_2}$ 约合 \qty{2.7e11}{t} 的
氧气。

\textbf{19.} $\num{1.2e15}/\num{2.7e11} \approx 4500$ 年的产量
（计入海洋的份额，约 2000 年）。但呼吸作用消耗氧气的速率
与\omterm{def:g10:cell-metabolism:autotrophy}{光合作用}产生氧气的速率几乎相同，
因此氧气储量是平衡而非累积；空气中的氧气
是数亿年间，由产量超过消耗的微小盈余
逐步积累而成。

\textbf{20.} 约 2\% 的全部日光（可用波长的 5\%）被储存为糖；
释放的氧气即被分解水的氧；
大气中约四分之一的碳每年
都经过\omterm{def:g10:cell-metabolism:autotrophy}{光合作用}。
\end{solution}
